\documentclass[a4paper]{article}

\usepackage{graphicx}
\usepackage{amsmath,amssymb}
\usepackage{minted}
\usepackage{multirow}
\usepackage{float}
\usepackage{todonotes}
\usepackage{forest}
\usepackage{xstring}
\usepackage{xcolor}
\usepackage[most]{tcolorbox}
\usepackage{xparse}
\usepackage{natbib}

\usetikzlibrary{graphs,graphdrawing}
\usegdlibrary{layered}

% for external graphviz
\input{/home/donkarlo/Dropbox/repo/nd_ai_project/src/nd_ai/knoweldge_representation/language/formal/computer/programming/tex/latex/code_clips/external_graphviz.tex}

% for tags
\input{/home/donkarlo/Dropbox/repo/nd_ai_project/src/nd_ai/knoweldge_representation/language/formal/computer/programming/tex/latex/parts/control_sequence/kind/semtag/semtag.tex}

% for links
\input{/home/donkarlo/Dropbox/repo/nd_ai_project/src/nd_ai/knoweldge_representation/language/formal/computer/programming/tex/latex/code_clips/seehere.tex}

\title{Ventral attention system}
\date{}

\begin{document}
    \maketitle
    \cite{scalf-2014-trial-history-effects-in-the-ventral-attentional-network}  investigates whether the ventral attentional network preserves information from previous trials and is affected when previously ignored stimuli become task-relevant. Using fMRI and the distractor preview effect, the authors found stronger activation in VAN-related regions, especially the right middle frontal gyrus and right supramarginal gyrus, when current targets had been rejected in the previous trial. The study argues that VAN is not only a stimulus-driven reorienting system, but also carries trial-history information that shapes current attentional selection.
    \bibliographystyle{plainnat}
    \bibliography{/home/donkarlo/Dropbox/repo/nd_shared_research/bibliography/bibliography.bib}
\end{document}